\subsection{Exclusion of the contracted horizontal section}

\begin{statement}{P4.54}{Proved}
In the remaining positive ruled support of Part~4.2,
the horizontal pole section is not contracted by the
normalization resolution. It consequently maps with
degree two or three onto a rational normalized conductor
component whose generic transverse germ is $D_4$.
No transverse-type restriction is imposed on the other
conductor component.
\end{statement}
\begin{proof}
By P4.53, a contracted section would be the sole smooth
elliptic exceptional curve $H$ of an isolated normal
hypersurface singularity on the first ambient strict
transform. Its self-intersection is $-d<0$, and the
actual exceptional coordinate $r$ has order one along
$H$. Write $\pi:\widetilde V\to(V,0)$ for this local
resolution. We establish the following three coordinate
weights and leading degrees, retaining the actual $r$:
\[
 (1,1,1;3),\qquad(1,1,2;4),\qquad(1,2,3;6).
                                                        \tag{4.54w.1}
\]

\paragraph{1. Analytic formal functions and the exact associated graded ring}

Put \(L=\mathcal O_H(-H)\), of degree \(d\), and
\[
 I_m=(\pi_*\mathcal O_{\widetilde V}(-mH))_0
     =\{a\in\mathcal O_{V,0}:\operatorname{ord}_H(a)\ge m\}.
\]
For \(m\ge1\), the restrictions of
\(\mathcal O(-mH)\) to successive infinitesimal neighbourhoods
of \(H\) have graded factors \(L^{m+j}\), \(j\ge0\).
Each has zero first cohomology, by Serre duality on the elliptic
curve and its strictly positive degree. Their extensions also
have zero first cohomology.

The analytic formal-functions theorem is Grauert1960,
Section6, Hauptsatz IIa, printed p.59; Hauptsatz I on the same
page supplies coherence \cite[\S6, Hauptsätze I and IIa]{Grauert60}. The maximal-ideal pullback and the
ideal of \(H\) have the same radical. On a finite cover of
the compact exceptional curve, Noetherianity gives one uniform
power of the latter contained in the former. Their powers are
therefore cofinal. Applying the theorem with
\(\mathcal O(-mH)\) gives a zero completion for
\(R^1\pi_*\mathcal O(-mH)\). This coherent module is
supported at the origin, so it is zero. Consequently
\[
 I_m/I_{m+1}\simeq H^0(H,L^m)\qquad(m\ge1),             \tag{4.54w.2}
\]
through the actual restriction maps in
\[
 0\longrightarrow\mathcal O(-(m+1)H)
 \longrightarrow\mathcal O(-mH)
 \longrightarrow L^m\longrightarrow0.
\]
These maps respect multiplication. Since normality gives
\(\pi_*\mathcal O_{\widetilde V}=\mathcal O_V\), and a
function vanishes at the contracted point exactly when its
restriction to \(H\) vanishes, \(I_1=\mathfrak m\).
The degree-zero piece is \(\mathbb C\). Thus
\[
 \operatorname{gr}_{I}\mathcal O_{V,0}
   =\bigoplus_{m\ge0}H^0(H,L^m)                         \tag{4.54w.3}
\]
as graded algebras, not only as vector spaces.

The same formal-functions theorem in degree zero identifies
\[
 \widehat{\mathcal O}_{V,0}
      \simeq\varprojlim_m H^0(mH,\mathcal O_{mH})
      \simeq\varprojlim_m\mathcal O_{V,0}/I_m.          \tag{4.54w.4}
\]
The second equality uses the just-proved first-direct-image
vanishing. This supplies the exact completeness and separation
needed in the leading-equation argument below; no unproved equivalence between a
divisorial and an ordinary degree filtration is being used.

Because \(\mathfrak m^2\subset I_2\), (4.54w.2) gives a
surjection
\(\mathfrak m/\mathfrak m^2\to H^0(H,L)\).
The source has dimension three and the target dimension \(d\).
Hence \(d\le3\).

\paragraph{2. The three section rings and their actual elliptic curves}

For \(d=3\), Riemann--Roch shows that \(L\) separates
points and tangent vectors: subtracting any length-two
subscheme leaves a degree-one line bundle with one section.
The complete linear series therefore embeds \(H\) as a smooth
plane cubic. The restriction sequence on \(\mathbb P^2\),
using \(H^1(\mathbb P^2,\mathcal O(k))=0\), gives its whole
section ring as
\[
 \mathbb C[x,y,z]/(F_3),\qquad \deg(x,y,z)=(1,1,1),
\]
where \(F_3=0\) is that smooth cubic.

For \(d=2\), the complete pencil is base-point free and gives
a degree-two map \(p:H\to\mathbb P^1\), with
\(L=p^*\mathcal O(1)\). Trace splits
\(p_*\mathcal O_H=\mathcal O\oplus B^{-1}\).
Its Euler characteristic is zero, so \(\deg B=2\).
The algebra structure is therefore a double-cover equation
\(z^2=F_4(x,y)\), where the binary quartic has four distinct
zeros because the double cover is smooth. Projection formula
gives the entire section ring with generator weights
\((1,1,2)\) and that single degree-four relation.

For \(d=1\), write \(L=\mathcal O_H(p)\).
Riemann--Roch gives a degree-one section \(x\), one new
degree-two section \(y\), and one new degree-three section
\(z\). Dividing by powers of \(x\) gives functions of pole
orders two and three at \(p\). Their Weierstrass relation
has weighted degree six. Equivalently, after graded changes
only for describing this section ring, it is
\[
 z^2=y^3+A x^4y+B x^6,
 \qquad 4A^3+27B^2\ne0,
 \qquad\deg(x,y,z)=(1,2,3).
\]
The pole-order count shows that the monomials with at most one
factor \(z\), reduced by this relation, give every space
\(H^0(H,L^m)\). Thus this is the whole section ring. Its
weighted projective curve is the actual smooth \(H\).
The coordinate choice below does not require completing the
square or applying any of these optional graded simplifications.

\paragraph{3. Retaining the given coordinate using only permitted changes}

Start with analytic ambient coordinates \((r,u,v)\), retaining
the given \(r\). Their images form a basis of
\(\mathfrak m/\mathfrak m^2\). The map to
\(I_1/I_2=H^0(L)\) has the rank just computed.

When \(d=3\), it is an isomorphism, so any linear complement
to \(r\) supplies the other two valuation-one coordinates.

When \(d=2\), choose a linear coordinate \(s\) whose
leading section complements that of \(r\), and choose the
remaining independent linear coordinate \(v_0\) in the
kernel of the degree-one restriction. Thus
\(\operatorname{ord}_H(v_0)\ge2\). Its degree-two leading
section must be independent of \(r^2,rs,s^2\). Otherwise
all degree-two initial sections of analytic functions would
lie in that three-dimensional span: a Taylor expansion shows
that every term of ordinary order at least three has valuation
at least three, and the linear terms of valuation two are
accounted for by \(v_0\). This contradicts
\(h^0(L^2)=4\). Therefore \(v_0\) has valuation exactly
two and is the required indecomposable degree-two generator.
The coordinates \((r,s,v_0)\) were obtained by an invertible
linear change.

When \(d=1\), first choose two independent linear coordinates
\(u_0,v_0\) in the kernel of the degree-one restriction.
At least one linear combination, denoted \(s\), has a
degree-two initial section independent of \(r^2\), since
\(h^0(L^2)=2\). Choose the remaining independent linear
coordinate \(v_1\), and subtract a scalar multiple of \(s\)
so that its degree-two initial term is \(\alpha r^2\).
Set
\[
                         v_2=v_1-\alpha r^2.
\]
This is a triangular analytic coordinate change, and
\(\operatorname{ord}_H(v_2)\ge3\). Its degree-three
initial section must be independent of \(r^3,rs\).
Indeed, if it were absent or in that span, any Taylor series
of valuation three would have degree-three initial section
in that same span: the possible terms are precisely the
linear \(v_2\) term and the weighted-degree-three monomials
\(r^3,rs\); its valuation-one and valuation-two terms
must first vanish. This contradicts \(h^0(L^3)=3\).
Thus \(v_2\) has valuation exactly three and supplies the
indecomposable degree-three generator. The three ambient
coordinates remain analytically invertible, because only
an invertible linear change and the displayed quadratic
translation were performed.

This argument proves directly that the chosen analytic
coordinates have the prescribed orders and generate the
graded ring. It avoids a circular appeal to filtered
Nakayama or to a coordinate normalization that might move
the fixed \(r\).

\paragraph{4. The actual weighted leading equation}

Write the defining analytic equation in the coordinates just
constructed, with the weights in (4.54w.1). Its weighted initial
form is the single defining relation of the graded section
ring, up to a nonzero scalar. Here is an explicit proof of that
last assertion.

Let \(F_D\) be the graded relation from the section-ring calculation above, expressed
in the actual chosen initial generators. The analytic function
\(F_D(r,s,v)\) has valuation at least \(D+1\).
At each successive valuation, (4.54w.3) writes its initial
section as a weighted homogeneous polynomial in the three
generators. Subtract that polynomial. This produces a formal
series \(G=F_D+F_{D+1}+\cdots\) whose image in every
\(\mathcal O_{V,0}/I_m\) is zero. Since all generator
weights are at most three, the correction polynomials tend
to zero in ordinary power-series degree; their sum is a
well-defined formal power series. By (4.54w.4), its image in
\(\widehat{\mathcal O}_{V,0}\) is zero.

The completed hypersurface ideal is principal, generated by
the actual analytic equation \(f\). Thus \(G=h f\) for
a formal series \(h\). Taking weighted initial forms gives
\(F_D=\operatorname{in}(h)\operatorname{in}(f)\).
The polynomial \(F_D\) is irreducible because its weighted
projective curve is the smooth connected elliptic curve
\(H\). Also \(f\) is nonconstant and vanishes at the
origin. Hence \(\operatorname{in}(h)\) is a nonzero
constant, and \(\operatorname{in}(f)\) is a scalar multiple
of \(F_D\). This proves (4.54w.1) for the actual defining
equation, with all higher terms retained.

\paragraph{5. The exact descent to the original blowdown coordinates}

In the chart with original coordinates
\(x=r u\), \(y=r v\), every change above has the form
\[
 s=a u+b v+c r,\qquad
 v'=a'u+b'v+c'r-\alpha r^2,
 \qquad ab'-a'b\ne0.
\]
It descends to the actual holomorphic coordinate map
\[
 \begin{aligned}
 x'&=r s=a x+b y+c r^2,\\
 y'&=r v'=a'x+b'y+c'r^2-\alpha r^3,\\
 r'&=r.
 \end{aligned}                                        \tag{4.54w.5}
\]
Its Jacobian determinant at the origin is \(ab'-a'b\),
so the inverse-function theorem supplies an analytic inverse.
This is the required morphism preserving the original
parameter and blowdown structure.


\paragraph{Leading polynomials in the original coordinates.}
Apply this calculation to the simple elliptic point supplied
by P4.53. The original equation is $r^3f_1$, and its original
coordinates are $(x,y,z)=(rs,rv,r)$. The three pairs of
original weights and degrees are
\[
 (2,2,1;6),\qquad (2,3,1;7),\qquad (3,4,1;9).
\]
For a strict-transform monomial $r^as^bv^c$, multiplication
by $r^3$ gives $x^by^cz^{3+a-b-c}$. Its original weighted
degree is three plus its strict-transform weighted degree.
Analyticity of the original equation makes every
negative-power coefficient zero. Thus the actual weighted
leading term is a polynomial of the indicated degree.
On $z\ne0$, the ratios
\[
     \left(x/z^{1+w_s},\ y/z^{1+w_v}\right)
\]
identify its projective component with the strict-transform
elliptic relation on $r\ne0$. This is a birational map
of the actual residue-field curves; their smooth
normalizations are the same elliptic curve $H$.

Anisotropic scaling with these weights and the compact-cycle
flow in P4.52 transports every pairing of the leading
Milnor fibre into the actual local torus fibre. By P4.51,
all these pairings must vanish. We now contradict this
in each of the three cases.

\paragraph{Degree six.}
The leading polynomial is exactly
\[
                    P_6(x,y,z)=G_3(x,y,z^2),
\]
where $G_3(X,Y,Z)$ defines a smooth projective cubic.
The map $(x,y,z)\mapsto(x,y,z^2)$ gives a finite double
cover from $F=\{P_6=1\}$ to $\{G_3=1\}$, branched over
\[
                       B=\{G_3(x,y,0)=1\}.
\]
The homogeneous cubic fibre has Euler number nine: it
is a threefold unramified cover of the complement of
an elliptic cubic in $\mathbb P^2$, whose Euler number
is three. Let $r\in\{1,2,3\}$ be the number of distinct
zeros of the nonzero binary cubic $G_3(x,y,0)$ on
$\mathbb P^1$. The curve $B$ is a threefold cover of
that punctured projective line. Hence
\[
          \chi(B)=3(2-r),\qquad \chi(F)=18-\chi(B)=12+3r.
\]

Weighted monodromy has order six. Its only
nonidentity element with a fixed set is its cube,
which negates $z$ and fixes $x,y$; its fixed set is
$B$. The other four nonidentity elements have no
fixed point on $F$. We choose an invariant
semialgebraic compact core large enough for $F$
and $B$. The finite-order Lefschetz character of
$H_0-H_1+H_2$ is therefore $\chi(F)$ at the identity,
$\chi(B)$ at the cube, and zero elsewhere. Its
trivial-character multiplicity is
\[
                       \frac{\chi(F)+\chi(B)}6=3.
\]
Since $H_0$ has only the trivial character, the
nontrivial-character dimension in $H_2$ is at least
$\chi(F)-3=9+3r\ge12$.

Every nonzero critical orbit of $P_6$ lies on $z=0$
at a tangency of the cubic with $Z=0$. If the
intersection orders are $e_i$, then $\sum e_i=3$.
At a tangency, smoothness of $G_3$ makes its
$Z$-derivative a unit, and the exact transverse
germ is $u^{e_i}+z^2$, after holomorphic coordinates
and a unit. Consequently
\[
                       M=\sum(e_i-1)=3-r\le2.
\]
This also shows reducedness of $P_6$: a repeated
factor would give a two-dimensional critical set,
whereas the derivative calculation gives only
these critical orbits. The variation inequality
(4.52b) contradicts the lower bound twelve for
the nontrivial monodromy rank.

\paragraph{Weighted curve cohomology.}
We record the two curve facts needed for the last
cases, with their quotient comparisons. For
$\mathbb P(a,b,c)=\operatorname{Proj}\mathbb C[u,v,w]$
and a reduced equation of weighted degree $e$,
the exact graded sequence sheafifies to
\[
 0\longrightarrow\mathcal O_{\mathbb P}(-e)
 \longrightarrow\mathcal O_{\mathbb P}
 \longrightarrow\mathcal O_D\longrightarrow0.
\]
The three standard affine charts give its \v Cech
complex. The middle cohomology of every
$\mathcal O_{\mathbb P}(k)$ is zero. The top
cohomology of $\mathcal O_{\mathbb P}(-e)$ has
basis the Laurent monomials $u^{-i}v^{-j}w^{-k}$,
$i,j,k\ge1$, of weighted degree $-e$. Thus
\[
 h^1(\mathcal O_D)=
 \#\{\hbox{nonnegative weighted monomials of degree }
                         e-(a+b+c)\},
 \qquad h^0(\mathcal O_D)=1.                         \tag{4.54a}
\]
In particular the arithmetic genus is zero below
degree $a+b+c$, and is one in degrees $a+b+c$
and $a+b+c+1$ when there is exactly one variable
of weight one.

At a cyclic quotient vertex of prime order two
or three, with no invariant linear ambient
coordinate, a smooth coarse curve has a smooth
plane pullback. To prove this, its normalized
branches form one orbit. A single singular
branch of multiplicity $m\ge2$ has a faithful
normalization action of order $p$; every invariant
ambient function has order at least $2m>p$,
whereas a regular quotient requires a function
of normalization order $p$. For $p$ transitively
permuted branches the stabilizer is trivial,
and every invariant function has order at least
two, whereas a regular quotient requires order
one. Both alternatives are impossible. The
remaining single branch is smooth, and its
reduced equation has the character of a nonzero
linear ambient coordinate. These arguments
use the actual analytic invariant rings:
averaging makes the invariants functor exact,
so every quotient-curve function lifts to an
invariant ambient function.

\paragraph{Degree seven.}
Consider weights $(1,2,3)$ and degree seven.
Formula (4.54a) gives arithmetic genus zero for
curves of degree at most five and arithmetic
genus one in degrees six and seven. Thus a
positive-genus component has degree six or
seven, genus one, and multiplicity one. There
is at most one further degree-one factor;
the polynomial is reduced.

In degree seven the whole curve $D$ is
irreducible with arithmetic and normalization
genus one. Its total delta invariant is zero,
so it is a smooth coarse elliptic curve.
At the quotient vertices of orders two and
three, the preceding argument makes its plane
pullback smooth. Elsewhere the ambient plane
is smooth. Thus the affine cone has no
nonzero critical orbit and $M=0$.

In degree six write $D=E\cup\{z=0\}$, with
$z$ the weight-one coordinate. The curve $E$
is smooth coarse elliptic. It cannot contain
either quotient vertex: its degree-six
equation has trivial character there, while
both local-coordinate characters are
nontrivial, so a smooth lift would require
an impossible invariant linear equation.
The genus formula for $D$ gives total delta
one. Its two smooth components therefore
meet once transversely in the smooth part
of the ambient weighted plane. There is
one ordinary-double critical orbit and
$M=1$.

The weighted plane has Euler number three,
from its affine-cell stratification. In the
two cases $e(D)=0$ or one, so its complement
has Euler number three or two. Both quotient
vertices belong to $D$, since degree seven
is divisible by neither two nor three.
The weighted $\mathbb C^*$ action is free
off $D$, so $\{P_7=1\}$ is an unramified
sevenfold cover of that complement. The
free order-seven monodromy gives nontrivial
$H_2$ rank at least eighteen or twelve,
respectively. This contradicts (4.52b),
whose upper bound is zero or one.

\paragraph{Degree nine.}
For weights $(1,3,4)$, (4.54a) gives
arithmetic genus zero below degree eight
and one in degrees eight and nine.
A positive-genus component must have one
of these two degrees, and multiplicity one.

An elliptic degree-eight component is
impossible. It necessarily contains the
weight-three vertex, since eight is not
divisible by three. The local plane
characters are $1,1$ modulo three, whereas
its equation has character two. The
prime-order quotient argument would give
a smooth lift, whose equation must have
character one. This is a contradiction.

Thus the entire degree-nine curve $D$ is
irreducible and smooth coarse elliptic.
It avoids the weight-three vertex: its
equation there has character zero,
incompatible with a smooth plane lift.
At the weight-four vertex the plane
characters are $1,3$ and its equation
has character one. We check the lift
without assuming that four is prime.
The transitive orbit of normalized
branches has size one, two, or four.

For four branches the stabilizer is
trivial and every invariant maximal-ideal
function has order at least two on a
branch, contradicting regularity of the
quotient. For two branches the stabilizer
has order two. A singular branch gives
invariant orders at least four, exceeding
two. Two smooth branches would have a
quadratic initial equation, but the
quadratic characters of the ambient
action are only zero and two, never one.
For one branch the normalization action
is faithful of order four. Both ambient
coordinate orders are odd, since their
characters are one and three. A singular
branch therefore has multiplicity at
least three, and all invariant functions
have order at least six, exceeding four.
These contradictions leave a single
smooth branch.

The degree-nine cone is therefore smooth
away from its vertex. It is singular
at the vertex, because no variable has
weight nine. Its hypersurface critical
point is isolated, so the isolated-point
Morse transport of P4.52 gives a spherical
class of square $-2$ in its Milnor fibre.
This again contradicts the zero pairing
transported from the actual torus.

All three possibilities for the sole
simple elliptic exceptional curve have
now been excluded. Together with P4.53
this proves that the horizontal pole
section cannot be contracted. P4R.18
then gives the asserted generic $D_4$
image and degree two or three. The
other conductor component retains its
unrestricted transverse type.

The exact weighted cohomology counts,
order-six characters, quotient quadratic
characters and an explicit smooth cubic
are recomputed in
\texttt{code/part4/simple\_elliptic\_weighted.py}
and its executed outputs.

\end{proof}
